\section{製造物責任を四つの判定に分ける}
\label{theory:pl:framework}
製造物責任法は、製品の欠陥による生命・身体・財産の被害について、製造業者等の民事上の賠償責任を定める。事故が起きたという一事だけで責任が成立するわけではない。まず対象となる物と責任主体を定め、次に欠陥、損害、因果関係を確認し、最後に免責と期間制限を検討する。この順序なら、対象外という判断を安全対策の不要という判断へ飛躍させずに済む。\autocite{pl_statute,pl_qa}

\subsection{過失の証明と欠陥の証明}
\label{theory:pl:proof}
PL法では製造業者の故意又は過失の立証を責任成立の要件としない。ただし、欠陥、損害、両者の因果関係は別の要件である。「過失を証明する必要がない」から「欠陥さえあれば全損害を賠償する」と推論してはならない。H26問7の5とH29問8の2は、救済の変化を短く述べた文として読む一方、独立した法律説明では省略された要件を補って理解する。

\subsection{その製品だけの損害}
\label{theory:pl:selfdamage}
3条ただし書は、損害が当該製造物についてのみ生じた場合をPL法の賠償対象から除く。テレビが発火してテレビだけが壊れた場合と、住宅へ延焼した場合を対比する。前者でも契約等による請求の可能性は別にあり、保証期間が過ぎたという事情だけで全ての救済が消えるとはいえない。自動二輪車の故障も、他の人又は物に被害があるかを先に確かめる。\autocite{pl_statute,pl_qa}

\section{有体物・動産・製造加工の三条件}
\label{theory:pl:scope}
日本法の対象は製造又は加工された動産である。電気、音響、単体ソフトウェア、サービスは有体物という条件を満たさず、土地や家屋自体は動産ではない。ガスやガソリンは有体物であるため、電気と同じ分類に置かない。水道についても、供給サービスと製造加工された物を区別する必要がある。\autocite{pl_definitions2018,pl_qa}

\subsection{単体ソフトウェアと組込み製品}
\label{theory:pl:software}
ソフトウェア自体が対象外でも、その不具合がロボットという製造物の安全性の欠如として現れ、負傷との因果関係が認められるなら、ロボットの製造業者等の責任が問題となる。H24問7のプリインストールされたOSはソフトウェアそのものを、同問のロボットは製品全体を指す。情報の欠陥を有体製品の欠陥へ結び付ける段階が解法の要点である。

\subsection{不動産に取り付けた動産}
\label{theory:pl:attachment}
製造物の引渡し後に建物へ取り付けたという理由だけで、製造時の欠陥の責任が消えるわけではない。エスカレータ、シャッター、サッシ等について、引渡時の動産性と欠陥、損害との因果関係を検討する。建物本体の施工不良とは異なる。逐条解説はエスカレータからの転落事件の裁判例も紹介している。\autocite{pl_definitions2018}

\subsection{冷凍・冷蔵と加工の区別}
\label{theory:pl:processing}
加工は、物の本質を保ちながら新しい属性を付け加える操作として理解する。単なる切断、冷凍、冷蔵、乾燥は基本的にこの加工へ含めない。一方、加熱、調味、塩漬け、燻製等により新しい属性を加えた食品は対象になり得る。冷凍肉とハムを比べる際には、保存状態の変化と新属性の付与を分ける。植物工場のトマトや養殖魚も、生産技術が高度であるという事情だけでは製造加工された物とはいえない。\autocite{pl_definitions2018}

\paragraph{初回解答を訂正する根拠}
H25問7の初回独立解答は、冷凍冷蔵を加工と扱って不適切6件・候補5とした。逐条解説の物理3頁（印刷4頁）を画像で確認し、この分類を訂正した。オを適切とすると、不適切はア・ウ・エ・カ・キの5件となり候補4に対応する。初回記録は保存し、公式解答との一致を訂正理由の代わりにしない。同問イの「品質上の欠陥」は安全性に関係する範囲で読む必要がある。

\section{欠陥は安全性の不足である}
\label{theory:pl:defect}
欠陥の判断には、製品の特性、通常予見される使用形態、引渡時期等を総合して考慮する。単なる性能の不足や外観の不良と、安全性の不足を区別する。設計、製造、指示警告の各段階に安全性の不足が生じ得る。メーカーが意図した操作だけを検討すればよいわけではなく、通常予見できる誤使用も安全設計の検討に含める。\autocite{pl_qa,pl_statute}

\section{製造業者等と免責の条件}
\label{theory:pl:producer}
業として製造・加工・輸入した者に加え、製造業者としての表示、製造業者と誤認させる表示、実質的製造業者と認められる表示をした者が責任主体となり得る。OEMの販売者は「常に除外」でも「常に該当」でもなく、表示と実質を確認する。単なる販売店と輸入者を一括して対象外としてはならない。\autocite{pl_statute,pl_qa}

\subsection{開発危険の抗弁}
\label{theory:pl:development}
引渡時の科学技術の知見では欠陥を認識できなかったことを製造業者等が証明した場合には、4条の免責が問題となる。担当者個人が知らなかったという説明では足りない。「認識できなくても必ず責任を負う」は免責規定を落とした文である。部品又は原材料には、欠陥が専ら完成品側の設計指示によって生じ、かつ部品側に過失がない場合の別の免責もある。

\section{修理された製品と修理行為を分ける}
\label{theory:pl:repair}
性質の回復・維持を行う修理、修繕、整備というサービス自体は、基本的に製造加工ではない。しかし、修理を経た元の製品が製造物でなくなるわけではない。元の引渡時にあった欠陥か、後の修理で生じた不具合かを区別する。H24問7の「使用者が修理した製品」を対象物の定義として読むと対象性は残り、修理行為や修理起因の不具合について読む場合とは数え方が異なる。この問の初回候補3と公式候補4の差は、解釈疑義として保存する。\autocite{pl_qa,pl_definitions2018}

\subsection{再生品と元の製造業者}
\label{theory:pl:recycle}
使用困難になった物の一部から製造加工した再生品は対象となる。再生した者が責任主体となり得るほか、原材料となった元の製品の引渡時の欠陥が再生品の損害と因果関係を持つなら、元の製造業者の責任も検討する。「最後に再生した者が全ての責任を負う」と一律にはいえない。\autocite{pl_definitions2018}

\section{期間制限は出題時点を確認する}
\label{theory:pl:time}
現行法の短期期間は損害及び賠償義務者を知った時から原則3年、生命又は身体侵害では5年である。長期期間は製造業者等の引渡しから10年であり、消費者の購入日から数えるとは限らない。蓄積性物質や潜伏期間のある損害には起算点の特則もある。H30問9の通常のテレビ火災について、欠陥の証明だけで期間経過を無視できるという文は成立しない。\autocite{pl_statute,pl_qa}

\subsection{2020年4月1日の改正と経過措置}
\label{theory:pl:transition}
改正前は生命身体侵害にも短期3年を用いた。2020年4月1日施行の改正ではこれを5年へ延長し、施行時に旧短期期間が完成していない場合の適用を定めた。長期10年についても時効期間であることを明記した。消費者庁の2018年逐条解説は原則として改正後の規定を説明する資料であるため、刊行年をそのまま当時の施行法の年とみなさない。\autocite{pl_transition2020,pl_time2018}

\section{民法・特別法・海外制度との関係}
\label{theory:pl:otherlaw}
PL法に特別の定めがない責任については民法の規定を用いる。対象外となる事例に全く責任がないとは結論できない。原子炉運転等による原子力損害については原子力損害賠償法の規定を確認し、同法4条3項によるPL法不適用を扱う。また、日本法だけを守れば海外の製造物責任制度も満たすという国際共通化の主張は成立しない。逐条解説にも電気の対象性の違いが説明されている。\autocite{pl_nuclear_statute,pl_statute,pl_definitions2018}

\section{賠償責任の回避を超える製品安全}
\label{theory:pl:csr}
R02問12は事故後の賠償だけを問う問題ではない。法令順守後のリスク低減、消費者の期待を踏まえた基準、事故予防、対話、弱者と将来への配慮、緊急時の被害拡大防止、苦情解決の手段を順に検討する。問題に提示された社会的責任の定義に照らすと、7項目はいずれも安全・安心へ向けた行動であり、不適切0件に対応する。ISO26000原規格全文との照合は未完として残す。\autocite{pl_product_csr2016}
